\chapter{The Buy-Side Trading Desk and Its Brokers}\label{ch:mx:the-buy-side-trading-desk-and-its-brokers}

A fund's trading desk sends its orders to six brokers' algorithms and pays each in commissions. Which broker is best cannot be seen in a quarter's
data, so the desk lets a wheel choose, and measures. This chapter describes what the desk does and the systems it does it with, then simulates a
year and more of its flow to answer three questions: how many months of orders it takes to tell the best broker from the second, what a
performance-weighted wheel gains and loses, and what happens to a broker ranking when the traders choose who gets the hard orders.

\section{What the desk does}

A portfolio manager decides what to buy and sell; the trading desk decides how. It receives the orders, chooses for each a broker, an algorithm and
its parameters (or a block negotiation), watches them while they work, and answers for the result under the firm's best-execution duty (One Quant Book
1, chapter~11). Two systems carry the work.

\begin{definition}[Order management system, execution management system]\label{def:mx:the-buy-side-trading-desk-and-its-brokers:oms}
An \emph{order management system}\index{order management system} (OMS) holds the fund's orders and positions from the portfolio manager's decision to
settlement: compliance checks, allocations to accounts, the record of what was ordered and done. An \emph{execution management system}\index{execution
management system} (EMS) is the trader's screen for working orders: market data, brokers' algorithms and direct market access, and the \omterm{def:mx:the-almgren-chriss-framework:parent}{child orders} and
fills in real time.
\end{definition}

The OMS knows the decision time and price; the EMS knows every \omterm{def:mx:the-almgren-chriss-framework:parent}{child order}. TCA (chapter~19) needs both, and the most common flaw in a desk's data is
that the decision time never left the portfolio manager's head.

\section{High touch and low touch}

\begin{definition}[High-touch and low-touch trading]\label{def:mx:the-buy-side-trading-desk-and-its-brokers:touch}
\emph{High-touch trading}\index{high-touch trading} hands an order to a person at the broker (a sales-trader or cash trader, One Quant Book 1, chapter~2),
who works it with judgement, finds natural counterparties for blocks, and may commit the broker's capital. \emph{Low-touch
trading}\index{low-touch trading} sends it to the broker's electronic channels (algorithms, direct market access, smart routing), with the buy-side
trader choosing and supervising.
\end{definition}

High touch is for orders whose size or information makes the market's reaction the main cost: a large block in an illiquid stock, a trade the fund
does not want seen (chapter~9's indications of interest and \omterm{def:mx:dark-pools-and-blocks:conditional}{conditional orders}). Low touch is for the rest, which is most orders by count. The line
moves with the order's difficulty, and so does the comparison of brokers: a broker that gets the hard orders looks expensive whatever it does.

\section{Algo wheels}

\begin{definition}[Algo wheel]\label{def:mx:the-buy-side-trading-desk-and-its-brokers:wheel}
An \emph{algo wheel}\index{algo wheel} allocates the desk's low-touch orders among brokers' algorithms by a rule rather than by the trader's choice:
at random, in rotation, stratified by the order's difficulty, or weighted by past performance, so that each broker's results can be compared on
comparable flow.
\end{definition}

A wheel is an A/B test (One Quant Book 7, chapter~21) whose randomisation unit is the order. The simulated desk of this chapter sends 600 orders a
month to six brokers. Each order's difficulty is its \omterm{def:mx:transaction-cost-analysis:pretrade}{pre-trade cost estimate} (chapter~19), lognormal with a mean of 15 basis points and a standard
deviation of 25; its cost is the broker's effect plus that difficulty plus noise, Student $t$ with four degrees of freedom and a standard deviation of 30
basis points (the heavy tails of real execution costs). The brokers' effects are 0, 3, 4, 5, 6 and 8 basis points: broker~1 is the best, broker~2 is 3
basis points worse.

The minimum detectable effect (One Quant Book 7, chapter~21) sets how long the wheel must turn. To detect a difference $\delta$ between two brokers with
a two-sided test at 5\% and a power of 80\% (One Quant Book 4, chapter~12), each needs $2\sigma^2(z_{0.975}+z_{0.8})^2/\delta^2$ orders; with six
brokers sharing 600 orders a month, $\delta=3$ and $\sigma=30$, that is 15.7 months. Without difficulty adjustment the noise includes the orders'
difficulty, $\sigma=\sqrt{30^2+25^2}=39.1$, and it takes 26.6 months.

\omcode{code/firm/algowheel/firm_algowheel.py}{78}{82}{The months of a uniform wheel needed to detect a difference between two brokers: the orders each broker needs for a two-sided test at the chosen power, times the number of brokers, over the desk's monthly flow.}\label{lst:mx:the-buy-side-trading-desk-and-its-brokers:months}

\Cref{fig:mx:the-buy-side-trading-desk-and-its-brokers:power} checks the formula by simulation: 300 desks, each testing broker~2 against broker~1
every two months. The adjusted test reaches 80\% power after 16 months, the raw one after 24. Stratifying the wheel by difficulty (each broker gets the
same mix of orders within twenty difficulty buckets) makes the raw comparison slightly more precise, a standard deviation of 1.51 basis points after a
year against 1.59 at random, but adjustment does far more: it removes the difficulty from the noise instead of balancing it.

\begin{omfigure}
\begin{tikzpicture}
\begin{axis}[width=9.6cm,height=5.6cm,xlabel={\small months of flow (600 orders a month)},ylabel={\small power (\%)},
  tick label style={font=\footnotesize},grid=major,grid style={gray!25},xmin=0,xmax=41,ymin=0,ymax=100,
  legend columns=2,legend style={font=\footnotesize,at={(0.5,-0.3)},anchor=north,draw=none,fill=none,
  /tikz/every even column/.append style={column sep=8pt}},table/col sep=comma]
\addplot[omDef,very thick,mark=*,mark size=1.4pt] table[x=months,y=adjusted]{figdata/microstructure/20-the-buy-side-trading-desk-and-its-brokers/power.csv};
\addplot[omProp,very thick,mark=square*,mark size=1.4pt] table[x=months,y=raw]{figdata/microstructure/20-the-buy-side-trading-desk-and-its-brokers/power.csv};
\addplot[gray,dashed,domain=0:41] {80};
\legend{adjusted for difficulty,raw}
\end{axis}
\end{tikzpicture}

\omcaption{The probability that a uniform wheel shows broker~2 significantly worse than broker~1 (a true difference of 3 basis points), against the
months of flow: 300 simulated desks. Data: \texttt{mx\_wheel.power\_study}.}\label{fig:mx:the-buy-side-trading-desk-and-its-brokers:power}
\end{omfigure}

A uniform wheel pays for its knowledge: it sends five-sixths of the flow to brokers that are not the best, 4.33 basis points an order above always
using broker~1. A performance-weighted wheel spends less. Thompson sampling keeps a posterior for each broker's adjusted cost and sends each order to
the broker whose sampled cost is lowest; it learns as it goes (\cref{fig:mx:the-buy-side-trading-desk-and-its-brokers:thompson}). Over 24 months it
sent 87.2\% of the orders to broker~1 (97.5\% in the last month) at an excess cost of 0.54 basis points an order. The price is information: broker~2
received 5.4\% of the flow, and the test of broker~1 against broker~2 had a power of 88\% after 24 months, against 100\% for the uniform wheel, which
also knows the other four brokers equally well. A wheel that must also keep brokers engaged and detect a broker that gets worse needs a floor on each
broker's share.

\begin{omfigure}
\begin{tikzpicture}
\begin{axis}[width=9.6cm,height=5.4cm,xlabel={\small month},ylabel={\small share of orders to broker 1 (\%)},
  tick label style={font=\footnotesize},grid=major,grid style={gray!25},xmin=0,xmax=25,ymin=0,ymax=100,
  legend columns=2,legend style={font=\footnotesize,at={(0.5,-0.3)},anchor=north,draw=none,fill=none,
  /tikz/every even column/.append style={column sep=8pt}},table/col sep=comma]
\addplot[omDef,very thick,mark=*,mark size=1.2pt] table[x=month,y=thompson]{figdata/microstructure/20-the-buy-side-trading-desk-and-its-brokers/thompson.csv};
\addplot[omProp,very thick] table[x=month,y=uniform]{figdata/microstructure/20-the-buy-side-trading-desk-and-its-brokers/thompson.csv};
\legend{Thompson sampling,uniform}
\end{axis}
\end{tikzpicture}

\omcaption{The share of each month's orders sent to the best broker by a Thompson-sampling wheel and by a uniform one, averaged over 60 simulated desks.
Data: \texttt{mx\_wheel.thompson\_study}.}\label{fig:mx:the-buy-side-trading-desk-and-its-brokers:thompson}
\end{omfigure}

\section{Evaluating brokers}

\begin{definition}[Broker scorecard]\label{def:mx:the-buy-side-trading-desk-and-its-brokers:scorecard}
A \emph{broker scorecard}\index{broker scorecard} is the desk's periodic report on each broker: orders and value traded, costs against the agreed
benchmarks raw and adjusted for difficulty with their standard errors, and softer measures (service, capital commitment, the quality of block
liquidity), from which the desk sets next period's allocation.
\end{definition}

\begin{table}[ht]
\centering
\small
\begin{tabular}{@{}lrrrrr@{}}
\toprule
broker & orders & raw cost & adjusted cost & s.e. & rank \\
\midrule
broker 1 & 1\,175 & 15.4 & 15.1 & 0.9 & 1 \\
broker 2 & 1\,246 & 16.7 & 17.7 & 0.9 & 2 \\
broker 3 & 1\,171 & 18.6 & 18.7 & 0.9 & 3 \\
broker 4 & 1\,169 & 21.9 & 21.0 & 0.9 & 5 \\
broker 5 & 1\,207 & 20.5 & 20.6 & 0.8 & 4 \\
broker 6 & 1\,232 & 22.9 & 22.7 & 0.8 & 6 \\
\bottomrule
\end{tabular}
\caption{A year of the uniform wheel as a scorecard: costs in basis points at the average difficulty, with robust standard errors. The true effects are
0, 3, 4, 5, 6 and 8 basis points above broker~1. Data: \texttt{mx\_wheel.card}.}\label{tab:mx:the-buy-side-trading-desk-and-its-brokers:card}
\end{table}

The scorecard of \cref{tab:mx:the-buy-side-trading-desk-and-its-brokers:card} ranks brokers 4 and 5 the wrong way round: their true difference is 1
basis point and the standard errors are 0.8 to 0.9. A ranking is a point estimate; the scorecard must show the errors. Anand, Irvine, Puckett and
Venkataraman (2012) found that some brokers deliver better executions consistently over time, so persistent differences exist, but finding them takes
years of a desk's flow or the pooled flow of many desks.

Without a wheel the comparison can be worse than noisy. Suppose the desk's traders send the hardest fifth of the orders to the broker they trust most,
broker~1, and the rest to the other five. After a year, the raw scorecard ranks broker~1 last, 35.1 basis points worse than broker~2; adjusted for
difficulty it ranks it first, 3.8 basis points better than the next. The adjustment works here only because the difficulty model is right; with a
wrong model, only randomisation protects the comparison.

\section{Paying for research and execution}

\begin{definition}[Commission sharing agreement, research unbundling]\label{def:mx:the-buy-side-trading-desk-and-its-brokers:csa}
A \emph{commission sharing agreement}\index{commission sharing agreement} lets a fund pay commissions to an executing broker and direct part of them to
other firms that provide research. \emph{Research unbundling}\index{research unbundling} separates the payment for research from the payment for
execution, so that each is priced and chosen on its own.
\end{definition}

Bundled commissions tie the choice of broker to the research the fund wants, which is another reason the desk's allocation may not follow execution
quality alone. The rules have moved several times.

\begin{dated}{2026-09}{Paying for research}\label{dat:mx:the-buy-side-trading-desk-and-its-brokers:research}
MiFID II, applied from 3 January 2018, required investment firms to separate charges for research from charges for execution: research paid from the
firm's own resources or from a research payment account agreed with clients. In the United Kingdom the FCA allowed joint payments for research and
execution again from 1 August 2024 (policy statement PS24/9), under conditions. In the EU, Directive (EU) 2024/2811 of 23 October 2024 (part of the
Listing Act) lets firms choose to pay jointly or separately, if they tell clients and assess the research each year; member states had to apply it
from 6 June 2026. In the United States, SEC staff no-action relief of 26 October 2017 let broker-dealers receive separate research payments from
MiFID II firms without being regulated as investment advisers; extended in 2019, it expired on 3 July 2023.
\end{dated}

For the desk, unbundling has a practical consequence: once research is paid separately, the execution budget can follow the scorecard.

\section{Tutorial: an algo wheel}\label{tut:mx:the-buy-side-trading-desk-and-its-brokers:tut}

\begin{tutorial}
\textbf{Goal.} Simulate a desk's flow over six brokers, run uniform, stratified and Thompson-sampling wheels, evaluate the brokers with and without
difficulty adjustment, and compute the months needed to rank them.
\textbf{End state:} \cref{fig:mx:the-buy-side-trading-desk-and-its-brokers:power,fig:mx:the-buy-side-trading-desk-and-its-brokers:thompson},
\cref{tab:mx:the-buy-side-trading-desk-and-its-brokers:card} and the numbers of sections 3 and 4.
\begin{tutsteps}
\item \textbf{Flow.} \texttt{mx\_wheel.desk(months, seed)}, \texttt{cost\_of(broker, flow)}.
\item \textbf{Allocate.} \texttt{firm\_algowheel.random\_allocation}, \texttt{stratified\_allocation}, \texttt{Thompson}.
\item \textbf{Evaluate.} \texttt{evaluate(cost, broker, difficulty, k, adjust)} on \texttt{firm.tca}'s clustered regression, \texttt{scorecard};
\texttt{months\_needed}.
\item \textbf{Studies.} \texttt{power\_study()}, \texttt{stratified\_study()}, \texttt{thompson\_study()}, \texttt{routing\_study()}, \texttt{card()};
draw with \texttt{fig\_wheel.py}.
\end{tutsteps}
\textbf{What to change next.} Let a broker's effect drift over time and see how fast each wheel notices; add a floor to Thompson sampling's shares;
use CUPED (One Quant Book 7, chapter~21) with the pre-trade estimate as the covariate.
\end{tutorial}

\section{Build: algo wheel}\label{bld:mx:the-buy-side-trading-desk-and-its-brokers:algowheel}

\begin{build}
\bfield{Purpose} The desk's allocation and evaluation of brokers, used by chapter~28's \omterm{def:mx:benchmark-algorithms:algo}{execution algorithm} to compare its own versions.
\bfield{Interface} \texttt{random\_allocation(n, k, rng)}, \texttt{stratified\_allocation(difficulty, k, buckets, rng)}, \texttt{Thompson(k,
prior\_mean, prior\_sd, noise\_sd)} with \texttt{choose} and \texttt{update}, \texttt{evaluate(cost, broker, difficulty, k, adjust, clusters)},
\texttt{months\_needed(delta, sd, orders\_per\_month, k, alpha, power)}, \texttt{scorecard(cost, broker, difficulty, k, names, clusters)}.
\bfield{Rules} Costs in basis points, positive when paid; difficulty is the pre-trade estimate, known before allocation; broker effects reported at the
average difficulty with clustered standard errors.
\bfield{Acceptance tests} \texttt{code/firm/algowheel/tests/}: stratification balances difficulty better than chance; the adjusted evaluation recovers
planted effects that the raw one distorts; the months formula; Thompson sampling concentrates on the best broker.
\bfield{Stretch} Time-varying effects; share floors; pooling across desks.
\end{build}

\begin{omsources}
\item A.~Anand, P.~Irvine, A.~Puckett and K.~Venkataraman, ``Performance of institutional trading desks: an analysis of persistence in trading
costs'', \emph{Review of Financial Studies} 25(2), 2012.
\item FCA, PS24/9, ``Payment optionality for investment research'', 2024.
\item Directive (EU) 2024/2811 amending Directive 2014/65/EU, 23 October 2024.
\item SEC, press release 2017-200 (26 October 2017), and Commissioner M.~T.~Uyeda, statement on the expiration of the staff no-action letter on
MiFID II, 5 July 2023.
\end{omsources}

\section{Exercises}

\begin{exercise}[$\star$]\label{exo:mx:the-buy-side-trading-desk-and-its-brokers:1}
How many orders does each broker need for a 3-basis-point difference to be detected with 80\% power at 5\%, with a noise of 30 basis points?
\end{exercise}

\begin{exercise}[$\star$]\label{exo:mx:the-buy-side-trading-desk-and-its-brokers:2}
The desk drops to three brokers. How many months does the adjusted comparison of the best two now need?
\end{exercise}

\begin{exercise}[$\star$]\label{exo:mx:the-buy-side-trading-desk-and-its-brokers:3}
What is the excess cost per order of a uniform wheel over always choosing broker~1?
\end{exercise}

\begin{exercise}[$\star\star$]\label{exo:mx:the-buy-side-trading-desk-and-its-brokers:4}
Why does stratification help the raw comparison so little here?
\end{exercise}

\begin{exercise}[$\star\star$]\label{exo:mx:the-buy-side-trading-desk-and-its-brokers:5}
Thompson sampling saves 3.8 basis points an order. What does it give up, and when would a desk prefer the uniform wheel?
\end{exercise}

\begin{exercise}[$\star\star$]\label{exo:mx:the-buy-side-trading-desk-and-its-brokers:6}
Why did the raw scorecard rank broker~1 last when the traders sent it the hardest orders, and what assumption does the adjustment rely on?
\end{exercise}

\begin{exercise}[$\star\star\star$]\label{exo:mx:the-buy-side-trading-desk-and-its-brokers:7}
\emph{Coding.} With the noise's standard deviation halved (15 basis points), how many months do the adjusted and raw comparisons need by the formula?
\end{exercise}

\begin{exercise}[$\star\star\star$]\label{exo:mx:the-buy-side-trading-desk-and-its-brokers:8}
\emph{Find the flaw.} ``After a quarter of the wheel, broker~3 is 2 basis points cheaper than broker~2, so we move its flow up.''
\end{exercise}

\section{Problem: Six Brokers, One Wheel}

\begin{problem}[{Weekend problem --- six brokers, one wheel}]\label{pb:mx:the-buy-side-trading-desk-and-its-brokers:1}
A desk's head asks how long its new wheel must run before it can say which broker is best, and whether it should weight the wheel by performance.

\medskip\noindent\textbf{Part I --- The desk.}
\begin{enumerate}
\item What do the OMS and the EMS each hold, and why does TCA need both?
\item Distinguish high-touch and \omterm{def:mx:the-buy-side-trading-desk-and-its-brokers:touch}{low-touch trading} and say which orders go where.
\item Why does the difficulty of the orders a broker receives distort its scorecard?
\item Define an \omterm{def:mx:the-buy-side-trading-desk-and-its-brokers:wheel}{algo wheel} and its randomisation unit.
\end{enumerate}

\medskip\noindent\textbf{Part II --- How long.}
\begin{enumerate}[resume]
\item Describe the simulated flow and the brokers' effects.
\item Derive the orders each broker needs for a two-sided test at 5\% with 80\% power.
\item Why does the raw comparison need more months than the adjusted one?
\item \emph{State the named result}: the months of flow needed to tell the best broker from the second with 80\% power, with and without difficulty
adjustment.
\item What does stratification add?
\end{enumerate}

\medskip\noindent\textbf{Part III --- Weighting by performance.}
\begin{enumerate}[resume]
\item How does Thompson sampling allocate?
\item Give its share of flow to the best broker, its excess cost and its power, against the uniform wheel's.
\item What would you add to protect the desk against a broker that gets worse?
\item Read the scorecard: what can and cannot be concluded about brokers 4 and 5?
\end{enumerate}

\medskip\noindent\textbf{Part IV --- Judgement and money.}
\begin{enumerate}[resume]
\item What did the traders' routing do to the raw ranking, and what did adjustment recover?
\item What did Anand, Irvine, Puckett and Venkataraman find about persistence?
\item Define a \omterm{def:mx:the-buy-side-trading-desk-and-its-brokers:csa}{commission sharing agreement} and \omterm{def:mx:the-buy-side-trading-desk-and-its-brokers:csa}{research unbundling}.
\item Summarise where research payment rules stand in the UK, the EU and the US.
\item Why does unbundling matter to the wheel?
\item What would you tell the head of the desk?
\item In one sentence: why does a desk randomise?
\end{enumerate}
\end{problem}

\section{Interview questions}

\begin{interviewq}[$\star$ \iqroles{trader}]\label{iq:mx:the-buy-side-trading-desk-and-its-brokers:1}
When would you give an order to a sales-trader rather than an algorithm?
\end{interviewq}

\begin{interviewq}[$\star\star$ \iqroles{researcher}]\label{iq:mx:the-buy-side-trading-desk-and-its-brokers:2}
How would you design an \omterm{def:mx:the-buy-side-trading-desk-and-its-brokers:wheel}{algo wheel}'s evaluation so that a broker cannot game it?
\end{interviewq}

\begin{interviewq}[$\star\star$ \iqroles{researcher}]\label{iq:mx:the-buy-side-trading-desk-and-its-brokers:3}
How long must a wheel run to rank brokers 3 basis points apart, and what shortens it?
\end{interviewq}

\begin{interviewq}[$\star\star$ \iqroles{developer}]\label{iq:mx:the-buy-side-trading-desk-and-its-brokers:4}
What must the OMS and EMS record for the desk's TCA and wheel to work?
\end{interviewq}

\begin{interviewq}[$\star\star$ \iqroles{mle}]\label{iq:mx:the-buy-side-trading-desk-and-its-brokers:5}
Would you run the wheel as a bandit? What are the risks?
\end{interviewq}

\begin{interviewq}[$\star\star\star$ \iqroles{bank}]\label{iq:mx:the-buy-side-trading-desk-and-its-brokers:6}
You run a broker's algorithm business. A client's wheel ranks you fifth of six after one quarter. What do you do?
\end{interviewq}
